\section{数据治理：版权、隐私与质量控制}

数据产业一旦进入规模化，就绕不开治理问题。LLM 数据涉及版权、隐私和来源透明度；机器人数据涉及真实环境采集、用户场景、传感器记录和安全责任；仿真数据则涉及物理假设和评测偏差。没有治理，数据越多，风险越大。

\subsection{治理三角}

\noindent\begin{longtable}{p{0.22\textwidth}p{0.34\textwidth}p{0.35\textwidth}}
\toprule
维度 & 需要回答的问题 & 失败后果 \\
\midrule
版权与授权 & 数据是否有合法来源和使用范围？ & 训练数据争议、下游商业风险。 \\
隐私与安全 & 数据是否包含个人、位置、家庭、工厂或商业秘密？ & 泄露敏感信息，无法进入企业场景。 \\
质量与可追溯 & 数据从哪里来，如何清洗，如何影响模型？ & 模型失败无法定位，Recipe 不可复现。 \\
\bottomrule
\end{longtable}

\begin{warningbox}{数据治理不是法务尾项}
如果数据来源、清洗和授权不可追溯，模型能力越强，风险越大。机器人和企业场景尤其如此，因为数据往往包含真实空间、真实设备和真实用户行为。
\end{warningbox}

\subsection{质量控制的四个问题}

一份数据进入训练前，至少要问四个问题：它是否代表目标任务；它是否包含模型缺的能力；它是否有可靠标签或反馈；它是否会引入错误捷径。对于仿真数据，还要额外问：仿真里学到的策略能否迁移到真实世界。

\begin{importantbox}{质量控制是 Recipe 的前置条件}
没有质量控制，Recipe 只是玄学配比；有了质量控制，Recipe 才能变成可复现的工程方法。
\end{importantbox}

\subsection{本章小结}

数据治理决定数据能不能长期使用。版权、隐私、可追溯和质量控制，是数据产业从“采集”走向“基础设施”的门槛。

